% Auto-generated from raw.txt
% Usage:
%   \vocabentry{word}{/ipa/ (pos)}{definition}{example}{note}

\vocabentry{Linguistics}
{/lɪŋˈɡwɪstɪks/ (n)}
{The scientific study of language and its structure.}
{Linguistics explores how languages are formed, how they change over time, and how people use them.}
{Đây là một danh từ dùng để chỉ ngôn ngữ học, ngành nghiên cứu khoa học về ngôn ngữ.}

\vocabentry{Bilingual}
{/ˌbaɪˈlɪŋɡwəl/ (adj)}
{Speaking two languages fluently.}
{Growing up in a multicultural household, she became bilingual in English and Spanish at a very young age.}
{Đây là một tính từ dùng để chỉ khả năng nói lưu loát hai ngôn ngữ.}

\vocabentry{Communication}
{/kəˌmjuːnɪˈkeɪʃn/ (n)}
{The imparting or exchanging of information or news.}
{Effective communication is essential for maintaining strong personal and professional relationships.}
{Đây là một danh từ dùng để chỉ sự giao tiếp hoặc truyền thông tin.}

\vocabentry{Vocabulary}
{/vəˈkæbjəleri/ (n)}
{The body of words used in a particular language.}
{Reading extensively is one of the best ways to expand your academic vocabulary.}
{Đây là một danh từ dùng để chỉ vốn từ vựng của một ngôn ngữ hoặc một cá nhân.}

\vocabentry{Grammar}
{/ˈɡræmər/ (n)}
{The whole system and structure of a language or of languages in general.}
{Mastering grammar is crucial for writing clear and professional documents.}
{Đây là một danh từ dùng để chỉ ngữ pháp, hệ thống quy tắc của ngôn ngữ.}

\vocabentry{Dialect}
{/ˈdaɪəlekt/ (n)}
{A particular form of a language which is peculiar to a specific region or social group.}
{Although they speak the same language, the dialects of the north and south are quite different.}
{Đây là một danh từ dùng để chỉ tiếng địa phương hoặc phương ngữ.}

\vocabentry{Phonetics}
{/fəˈnetɪks/ (n)}
{The study and classification of speech sounds.}
{Phonetics helps language learners understand the correct way to produce sounds in a foreign language.}
{Đây là một danh từ dùng để chỉ ngữ âm học, nghiên cứu về các âm thanh trong lời nói.}

\vocabentry{Semantics}
{/sɪˈmæntɪks/ (n)}
{The branch of linguistics concerned with meaning.}
{Semantics is important because a single word can have multiple meanings depending on the context.}
{Đây là một danh từ dùng để chỉ ngữ nghĩa học, nghiên cứu về ý nghĩa của từ và câu.}

\vocabentry{Syntax}
{/ˈsɪntæks/ (n)}
{The arrangement of words and phrases to create well-formed sentences in a language.}
{English syntax typically follows a Subject-Verb-Object pattern.}
{Đây là một danh từ dùng để chỉ cú pháp, cách sắp xếp từ ngữ trong câu.}

\vocabentry{Eloquence}
{/ˈeləkwəns/ (n)}
{Fluent or persuasive speaking or writing.}
{The politician's eloquence won over many undecided voters during the televised debate.}
{Đây là một danh từ dùng để chỉ sự hùng biện hoặc khả năng diễn đạt lưu loát và cuốn hút.}

\vocabentry{Pragmatics}
{/præɡˈmætɪks/ (n)}
{The branch of linguistics dealing with language in use and the contexts in which it is used.}
{Pragmatics helps us understand the social rules of conversation, such as taking turns and being polite.}
{Đây là một danh từ dùng để chỉ ngữ dụng học, nghiên cứu cách sử dụng ngôn ngữ trong thực tế.}

\vocabentry{Fluency}
{/ˈfluːənsi/ (n)}
{The ability to speak or write a foreign language easily and accurately.}
{To achieve fluency in a new language, consistent practice and immersion are required.}
{Đây là một danh từ dùng để chỉ sự lưu loát hoặc trôi chảy trong ngôn ngữ.}

\vocabentry{Accent}
{/ˈæksent/ (n)}
{A distinctive mode of pronunciation of a language, especially one associated with a particular locality or social class.}
{He has a strong Scottish accent that can be difficult for outsiders to understand.}
{Đây là một danh từ dùng để chỉ giọng điệu hoặc trọng âm.}

\vocabentry{Translation}
{/trænzˈleɪʃn/ (n)}
{The process of translating words or text from one language into another; the product of such a process.}
{This book is a translation from the original Japanese text.}
{Đây là một danh từ dùng để chỉ sản phẩm dịch thuật, văn bản được chuyển từ ngôn ngữ này sang ngôn ngữ khác}

\vocabentry{Interpretation}
{/ɪnˌtɜːrprɪˈteɪʃn/ (n)}
{The action of explaining the meaning of something; or oral translation.}
{Simultaneous interpretation is a demanding task that requires high levels of concentration.}
{Đây là một danh từ dùng để chỉ sự giải thích hoặc thông dịch lời nói.}

\vocabentry{Colloquial}
{/kəˈloʊkwiəl/ (adj)}
{(Of language) used in ordinary or familiar conversation; not formal or literary.}
{It is important to learn colloquial expressions to understand how people actually speak in real life.}
{Đây là một tính từ dùng để chỉ ngôn ngữ thông tục, khẩu ngữ hằng ngày.}

\vocabentry{Jargon}
{/ˈdʒɑːrɡən/ (n)}
{Special words or expressions that are used by a particular profession or group and are difficult for others to understand.}
{Medical jargon can be very confusing for patients who do not have a background in medicine.}
{Đây là một danh từ dùng để chỉ thuật ngữ chuyên ngành.}

\vocabentry{Slang}
{/slæŋ/ (n)}
{A type of language that consists of words and phrases that are regarded as very informal.}
{Internet slang changes rapidly, with new words appearing and disappearing every year.}
{Đây là một danh từ dùng để chỉ tiếng lóng.}

\vocabentry{Literacy}
{/ˈlɪtərəsi/ (n)}
{The ability to read and write.}
{Improving adult literacy rates is a key goal for the government's education policy.}
{Đây là một danh từ dùng để chỉ trình độ học vấn hoặc khả năng biết đọc biết viết.}

\vocabentry{Nuance}
{/ˈnuːɑːns/ (n)}
{A subtle difference in or shade of meaning, expression, or sound.}
{A good translator must be able to capture the subtle nuances of the original text.}
{Đây là một danh từ dùng để chỉ sắc thái tinh tế trong ý nghĩa hoặc cách diễn đạt.}

\vocabentry{Discourse}
{/ˈdɪskɔːrs/ (n)}
{Written or spoken communication or debate.}
{Political discourse should be based on respect and a willingness to listen to different perspectives.}
{Đây là một danh từ dùng để chỉ diễn ngôn hoặc quá trình thảo luận.}

\vocabentry{Etymology}
{/ˌetɪˈmɑːlədʒi/ (n)}
{The study of the origin of words and the way in which their meanings have changed throughout history.}
{The etymology of many English words can be traced back to Latin or Greek.}
{Đây là một danh từ dùng để chỉ từ nguyên học.}

\vocabentry{Connotation}
{/ˌkɑːnəˈteɪʃn/ (n)}
{An idea or feeling that a word invokes in addition to its literal or primary meaning.}
{The word "cheap" has a negative connotation, suggesting poor quality as well as low price.}
{Đây là một danh từ dùng để chỉ nghĩa hàm ẩn hoặc cảm giác mà từ ngữ gợi ra.}

\vocabentry{Denotation}
{/ˌdiːnoʊˈteɪʃn/ (n)}
{The literal or primary meaning of a word, in contrast to the feelings or ideas that the word suggests.}
{While the denotation of "home" is a place where someone lives, its connotation is much warmer.}
{Đây là một danh từ dùng để chỉ nghĩa đen hoặc nghĩa từ điển.}

\vocabentry{Non-verbal}
{/ˌnɑːn ˈvɜːrbl/ (adj)}
{Not involving or using words or speech.}
{Non-verbal communication, such as body language, plays a huge role in how we perceive others.}
{Đây là một tính từ dùng để chỉ giao tiếp phi ngôn ngữ như cử chỉ, ánh mắt.}

\vocabentry{Articulation}
{/ɑːrˌtɪkjuˈleɪʃn/ (n)}
{The formation of clear and distinct sounds in speech.}
{Good articulation is essential for public speakers who want to be understood by a large audience.}
{Đây là một danh từ dùng để chỉ cách phát âm rõ ràng.}

\vocabentry{Cognate}
{/ˈkɑːɡneɪt/ (n)}
{(Of a word) having the same linguistic derivation as another.}
{"Night" in English and "Nuit" in French are cognates derived from the same Indo-European root.}
{Đây là một danh từ dùng để chỉ từ cùng gốc giữa các ngôn ngữ khác nhau.}

\vocabentry{Acquisition}
{/ˌækwɪˈzɪʃn/ (n)}
{The learning or developing of a skill, habit, or quality; in language, the natural process of learning.}
{Language acquisition in children happens naturally through exposure to their environment.}
{Đây là một danh từ dùng để chỉ sự thụ đắc ngôn ngữ.}

\vocabentry{Proverb}
{/ˈprɑːvɜːrb/ (n)}
{A short pithy saying in general use, stating a general truth or piece of advice.}
{"A picture is worth a thousand words" is a famous proverb that highlights the power of visual communication.}
{Đây là một danh từ dùng để chỉ tục ngữ hoặc câu ngạn ngữ.}

\vocabentry{Idiom}
{/ˈɪdiəm/ (n)}
{A group of words established by usage as having a meaning not deducible from those of the individual words.}
{"Breaking the ice" is a common idiom that means to start a conversation in a social setting.}
{Đây là một danh từ dùng để chỉ thành ngữ.}

\vocabentry{Euphemism}
{/ˈjuːfəmɪzəm/ (n)}
{A mild or indirect word or expression substituted for one considered to be too harsh or blunt.}
{"Passed away" is a common euphemism for the word "died."}
{Đây là một danh từ dùng để chỉ uyển ngữ hoặc lối nói giảm nói tránh.}

\vocabentry{Vernacular}
{/vərˈnækjələr/ (n)}
{The language or dialect spoken by the ordinary people in a particular country or region.}
{The poet chose to write in the vernacular to make his work more accessible to common people.}
{Đây là một danh từ dùng để chỉ tiếng mẹ đẻ hoặc ngôn ngữ bản xứ.}

\vocabentry{Monolingual}
{/ˌmɑːnəˈlɪŋɡwəl/ (adj)}
{(Of a person or society) speaking only one language.}
{In an increasingly globalized world, being monolingual can be a disadvantage in the job market.}
{Đây là một tính từ dùng để chỉ khả năng chỉ nói được một ngôn ngữ duy nhất.}

\vocabentry{Multilingual}
{/ˌmʌltiˈlɪŋɡwəl/ (adj)}
{Involving or using several languages.}
{Switzerland is a multilingual country with four official languages: German, French, Italian, and Romansh.}
{Đây là một tính từ dùng để chỉ tính đa ngôn ngữ.}

\vocabentry{Polyglot}
{/ˈpɑːliɡlɑːt/ (n)}
{A person who knows and is able to use several languages.}
{As a professional diplomat, he is a polyglot who speaks six different languages fluently.}
{Đây là một danh từ dùng để chỉ người biết nhiều thứ tiếng.}

\vocabentry{Homonym}
{/ˈhɑːmənɪm/ (n)}
{Each of two or more words having the same spelling or pronunciation but different meanings and origins.}
{The word "bank" is a homonym, meaning either a financial institution or the side of a river.}
{Đây là một danh từ dùng để chỉ từ đồng âm hoặc từ đồng hình.}

\vocabentry{Synonym}
{/ˈsɪnənɪm/ (n)}
{A word or phrase that means exactly or nearly the same as another word or phrase in the same language.}
{"Happy" and "joyful" are synonyms that can often be used interchangeably.}
{Đây là một danh từ dùng để chỉ từ đồng nghĩa.}

\vocabentry{Antonym}
{/ˈæntənɪm/ (n)}
{A word opposite in meaning to another.}
{"Hot" and "cold" are simple examples of antonyms.}
{Đây là một danh từ dùng để chỉ từ trái nghĩa.}

\vocabentry{Metaphor}
{/ˈmetəfɔːr/ (n)}
{A figure of speech in which a word or phrase is applied to an object or action to which it is not literally applicable.}
{The phrase "life is a roller coaster" is a metaphor for the ups and downs we all experience.}
{Đây là một danh từ dùng để chỉ phép ẩn dụ.}

\vocabentry{Simile}
{/ˈsɪməli/ (n)}
{A figure of speech involving the comparison of one thing with another thing of a different kind.}
{"As brave as a lion" is a common simile used to describe someone's courage.}
{Đây là một danh từ dùng để chỉ phép so sánh.}

\vocabentry{Hyperbole}
{/haɪˈpɜːrbəli/ (n)}
{Exaggerated statements or claims not meant to be taken literally.}
{When she said she was so hungry she could eat a horse, she was using hyperbole.}
{Đây là một danh từ dùng để chỉ phép ngoa dụ hoặc sự nói quá.}

\vocabentry{Irony}
{/ˈaɪrəni/ (n)}
{The expression of one's meaning by using language that normally signifies the opposite, typically for humorous or emphatic effect.}
{There was a sense of irony in the fact that the fire station burned down.}
{Đây là một danh từ dùng để chỉ sự mỉa mai hoặc trớ trêu.}

\vocabentry{Sarcasm}
{/ˈsɑːrkæzəm/ (n)}
{The use of irony to mock or convey contempt.}
{"Oh, brilliant!" he said with heavy sarcasm after I dropped the tray of glasses.}
{Đây là một danh từ dùng để chỉ sự châm biếm hoặc lời nói mỉa mai cay độc.}

\vocabentry{Ambiguity}
{/ˌæmbɪˈɡjuːəti/ (n)}
{The quality of being open to more than one interpretation; inexactness.}
{Avoid ambiguity in your legal documents to prevent future disputes.}
{Đây là một danh từ dùng để chỉ tính áp định hoặc mơ hồ, quan trọng trong giao tiếp rõ ràng}

\vocabentry{Redundancy}
{/rɪˈdʌndənsi/ (n)}
{The use of words or data that could be omitted without loss of meaning or function.}
{The phrase "I personally think" is often criticized for its redundancy.}
{Trong ngôn ngữ, đây là danh từ chỉ sự dư thừa lời văn.}

\vocabentry{Coherence}
{/koʊˈhɪrəns/ (n)}
{The quality of being logical and consistent.}
{To write a good essay, you must ensure that there is a clear coherence between your ideas.}
{Trong viết lách, đây là danh từ chỉ tính mạch lạc và nhất quán.}

\vocabentry{Cohesion}
{/koʊˈhiːʒn/ (n)}
{The action or fact of forming a united whole; in linguistics, the grammatical and lexical linking within a text.}
{Using transition words like "however" and "moreover" improves the cohesion of your writing.}
{Đây là một danh từ dùng để chỉ sự liên kết về mặt từ ngữ.}

\vocabentry{Cliché}
{/kliːˈʃeɪ/ (n)}
{A phrase or opinion that is overused and betrays a lack of original thought.}
{"Thinking outside the box" has become a bit of a cliché in the business world.}
{Đây là một danh từ dùng để chỉ lời nói rập khuôn hoặc cụm từ đã quá nhàm chán.}

\vocabentry{Mother tongue}
{/ˌmʌðər ˈtʌŋ/ (n)}
{The language which a person has grown up speaking from early childhood.}
{Although she speaks English fluently, her mother tongue is actually Vietnamese.}
{Đây là một cụm danh từ dùng để chỉ tiếng mẹ đẻ.}

\vocabentry{Lingua franca}
{/ˌlɪŋɡwə ˈfræŋkə/ (n)}
{A language that is adopted as a common language between speakers whose native languages are different.}
{English has become the global lingua franca for business, science, and technology.}
{Đây là một cụm danh từ dùng để chỉ ngôn ngữ chung/ngôn ngữ giao tiếp quốc tế.}

\vocabentry{Proficiency}
{/prəˈfɪʃnsi/ (n)}
{A high degree of skill; expertise.}
{Most universities require international students to demonstrate a high level of English proficiency.}
{Đây là một danh từ dùng để chỉ sự thông thạo hoặc trình độ thành thạo ngôn ngữ.}

\vocabentry{Impairment}
{/ɪmˈpermənt/ (n)}
{The state or fact of being impaired, especially in a specified faculty.}
{Sign language provides an essential tool for people with hearing impairment to communicate.}
{Trong giao tiếp, đây là sự khiếm khuyết, ví dụ "hearing impairment" - khiếm thính.}

\vocabentry{Sign language}
{/ˈsaɪn ˌlæŋɡwɪdʒ/ (n)}
{A system of communication using visual gestures and signs.}
{American Sign Language (ASL) is a complete and complex language with its own grammar rules.}
{Đây là một cụm danh từ dùng để chỉ ngôn ngữ ký hiệu.}

\vocabentry{Gesture}
{/ˈdʒestʃər/ (n/v)}
{A movement of part of the body, especially a hand or the head, to express an idea or meaning.}
{In some cultures, a nod of the head is a gesture that means "no" instead of "yes."}
{Đây là một danh từ/động từ chỉ cử chỉ hoặc điệu bộ.}

\vocabentry{Intonation}
{/ˌɪntəˈneɪʃn/ (n)}
{The rise and fall of the voice in speaking.}
{Changing your intonation can turn a simple statement into a question.}
{Đây là một danh từ dùng để chỉ ngữ điệu.}

\vocabentry{Phonology}
{/fəˈnɑːlədʒi/ (n)}
{The branch of linguistics that deals with systems of sounds.}
{Phonology examines how sounds function within a particular language system.}
{Đây là một danh từ dùng để chỉ âm vị học.}

\vocabentry{Morphology}
{/mɔːrˈfɑːlədʒi/ (n)}
{The study of the forms of words.}
{Morphology studies how prefixes and suffixes are added to root words to change their meaning.}
{Đây là một danh từ dùng để chỉ từ thái học/hình thái học, nghiên cứu cấu trúc của từ.}

\vocabentry{Lexicon}
{/ˈleksɪkən/ (n)}
{The vocabulary of a person, language, or branch of knowledge.}
{Each scientific discipline has its own unique lexicon of technical terms.}
{Đây là một danh từ dùng để chỉ từ vựng hoặc hệ thống từ vựng chuyên biệt.}

\vocabentry{Etiquette}
{/ˈetɪket/ (n)}
{The customary code of polite behavior in society or among members of a particular profession or group.}
{Understanding local etiquette is essential for effective cross-cultural communication.}
{Trong giao tiếp, đây là nghi thức xã giao.}

\vocabentry{Rhetoric}
{/ˈretərɪk/ (n)}
{The art of effective or persuasive speaking or writing.}
{The speaker used powerful rhetoric to inspire the audience to take action.}
{Đây là một danh từ dùng để chỉ tu từ học hoặc nghệ thuật hùng biện.}

\vocabentry{Interpreting}
{/ɪnˈtɜːrprɪtɪŋ/ (n)}
{The action of explaining the meaning of something; oral translation.}
{Interpreting at international conferences requires intense concentration and quick thinking.}
{Đây là một danh từ dùng để chỉ nghề thông dịch.}

\vocabentry{Transcription}
{/trænˈskrɪpʃn/ (n)}
{A written or printed representation of something.}
{The journalist made a full transcription of the recorded interview.}
{Đây là một danh từ dùng để chỉ sự phiên âm hoặc ghi chép lại lời nói.}

\vocabentry{Vocal}
{/ˈvoʊkl/ (adj)}
{Relating to the human voice.}
{Proper vocal technique is important for singers and public speakers alike.}
{Đây là một tính từ dùng để chỉ những gì thuộc về giọng nói.}

\vocabentry{Audible}
{/ˈɔːdəbl/ (adj)}
{Able to be heard.}
{His voice was barely audible over the noise of the traffic.}
{Đây là một tính từ dùng để chỉ âm thanh có thể nghe thấy được.}

\vocabentry{Inaudible}
{/ɪnˈɔːdəbl/ (adj)}
{Unable to be heard.}
{The dog whistle produces a sound that is inaudible to human ears.}
{Đây là một tính từ dùng để chỉ âm thanh không thể nghe thấy được.}

\vocabentry{Feedback}
{/ˈfiːdbæk/ (n)}
{Information about reactions to a product, a person's performance of a task, etc.}
{Constructive feedback from your teacher can help you improve your speaking skills.}
{Trong giao tiếp, đây là thông tin phản hồi.}

\vocabentry{Interaction}
{/ˌɪntərˈækʃn/ (n)}
{Reciprocal action or influence.}
{Social media has transformed the way we have social interactions with our friends.}
{Đây là một danh từ dùng để chỉ sự tương tác qua lại giữa người với người.}

\vocabentry{Intercultural}
{/ˌɪntərˈkʌltʃərəl/ (adj)}
{Relating to or involving different cultures.}
{The university offers a course on intercultural communication to prepare students for a global career.}
{Đây là một tính từ dùng để chỉ tính chất liên văn hóa.}

\vocabentry{Assimilation}
{/əˌsɪməˈleɪʃn/ (n)}
{The process of taking in and fully understanding information or ideas; in language, the merging of sounds.}
{The assimilation of new words from other languages is a common part of linguistic evolution.}
{Đây là một danh từ dùng để chỉ sự đồng hóa văn hóa hoặc âm thanh.}

\vocabentry{Acculturation}
{/əˌkʌltʃəˈreɪʃn/ (n)}
{The process of social, psychological, and cultural change that stems from the balancing of two cultures.}
{Acculturation often occurs when someone moves to a new country and starts learning the language.}
{Đây là một danh từ dùng để chỉ sự tiếp nhận văn hóa.}

\vocabentry{Linguistic}
{/lɪŋˈɡwɪstɪk/ (adj)}
{Relating to language or linguistics.}
{The country's linguistic diversity is one of its greatest strengths.}
{Đây là một tính từ dùng để chỉ những gì thuộc về ngôn ngữ học.}

\vocabentry{Sociolinguistics}
{/ˌsoʊsioʊlɪŋˈɡwɪstɪks/ (n)}
{The study of language in relation to social factors.}
{Sociolinguistics explores how factors like class, gender, and age affect the way we speak.}
{Đây là một danh từ dùng để chỉ xã hội ngôn ngữ học.}

\vocabentry{Psycholinguistics}
{/ˌsaɪkoʊlɪŋˈɡwɪstɪks/ (n)}
{The study of the relationships between linguistic behavior and psychological processes.}
{Psycholinguistics examines how the human brain produces and understands language.}
{Đây là một danh từ dùng để chỉ tâm lý ngôn ngữ học.}

\vocabentry{Bilingualism}
{/baɪˈlɪŋɡwəlɪzəm/ (n)}
{Fluency in or use of two languages.}
{Bilingualism has been shown to have several cognitive benefits for the brain.}
{Đây là một danh từ dùng để chỉ chủ nghĩa song ngữ hoặc tình trạng biết hai thứ tiếng.}

\vocabentry{Diglossia}
{/daɪˈɡlɑːsiə/ (n)}
{A situation in which two languages (or two varieties of the same language) are used under different conditions within a community.}
{In many Arab countries, there is a state of diglossia between the formal written language and the spoken dialects.}
{Đây là một danh từ chuyên ngành chỉ tình trạng song ngữ phân hóa chức năng.}

\vocabentry{Code-switching}
{/ˈkoʊd ˌswɪtʃɪŋ/ (n)}
{The practice of alternating between two or more languages or varieties of language in conversation.}
{Code-switching is common among bilingual people who move between different social groups.}
{Đây là một danh từ dùng để chỉ sự chuyển mã/đổi ngôn ngữ khi đang nói chuyện.}

\vocabentry{Pidgin}
{/ˈpɪdʒɪn/ (n)}
{A grammatically simplified means of communication that develops between two or more groups that do not have a language in common.}
{A pidgin language often develops in trading ports where people from many different countries meet.}
{Đây là một danh từ dùng để chỉ ngôn ngữ bồi.}

\vocabentry{Creole}
{/ˈkriːoʊl/ (n)}
{A mother tongue that originates from contact between two languages. (Đây là một danh từ dùng để chỉ ngôn ngữ hòa trộn (phát triển từ tiếng bồi thành tiếng mẹ đẻ).).}
{Haitian Creole is a vibrant language that developed from a mixture of French and West African languages.}
{}

\vocabentry{Orthography}
{/ɔːrˈθɑːɡrəfi/ (n)}
{The conventional spelling system of a language.}
{English orthography can be very challenging for learners due to its many irregular spellings.}
{Đây là một danh từ dùng để chỉ chính tả hoặc hệ thống chữ viết.}

\vocabentry{Script}
{/skrɪpt/ (n)}
{Handwriting as distinct from print; a particular system of writing.}
{The Arabic script is written from right to left and is used by millions of people.}
{Đây là một danh từ dùng để chỉ hệ thống chữ viết hoặc kịch bản.}

\vocabentry{Alphabet}
{/ˈælfəbet/ (n)}
{A set of letters or symbols in a fixed order, used to represent the basic sounds of a language.}
{The Greek alphabet was the first to include separate symbols for vowels.}
{Đây là một danh từ dùng để chỉ bảng chữ cái.}

\vocabentry{Logogram}
{/ˈlɔːɡəɡræm/ (n)}
{A sign or character representing a word or phrase.}
{Chinese characters are examples of logograms where each symbol represents a whole word.}
{Đây là một danh từ dùng để chỉ chữ biểu từ, ví dụ như chữ Hán.}

\vocabentry{Pictogram}
{/ˈpɪktəɡræm/ (n)}
{A pictorial symbol for a word or phrase.}
{Ancient Egyptian hieroglyphs were originally a system of pictograms.}
{Đây là một danh từ dùng để chỉ chữ tượng hình hoặc biểu đồ bằng hình ảnh.}

\vocabentry{Ideogram}
{/ˈɪdiəɡræm/ (n)}
{A written character symbolizing the idea of a thing without indicating the sounds used to say it.}
{Modern emojis can be seen as a type of universal ideogram used in digital communication.}
{Đây là một danh từ dùng để chỉ chữ biểu ý.}

\vocabentry{Punctuation}
{/ˌpʌŋktʃuˈeɪʃn/ (n)}
{The marks, such as period, comma, and parentheses, used in writing.}
{Correct punctuation is essential for ensuring that your writing is clear and easy to read.}
{Đây là một danh từ dùng để chỉ hệ thống dấu câu.}

\vocabentry{Inflection}
{/ɪnˈflekʃn/ (n)}
{A change in the form of a word (typically the ending) to express a grammatical function.}
{Latin is a highly inflected language where word endings change depending on their role in the sentence.}
{Đây là một danh từ dùng để chỉ sự biến cách hoặc sự biến đổi hình thái từ.}

\vocabentry{Conjugation}
{/ˌkɑːndʒuˈɡeɪʃn/ (n)}
{The variation of the form of a verb in an inflected language.}
{Learning the irregular conjugations of French verbs takes a lot of time and practice.}
{Đây là một danh từ dùng để chỉ sự chia động từ.}

\vocabentry{Declension}
{/dɪˈklenʃn/ (n)}
{The variation of the form of a noun, pronoun, or adjective.}
{German has four cases, and understanding the declension of articles is vital for learners.}
{Đây là một danh từ dùng để chỉ sự biến cách của danh từ, đại từ hoặc tính từ.}

\vocabentry{Register}
{/ˈredʒɪstər/ (n)}
{A variety of a language used for a particular purpose or in a particular social setting.}
{Using a formal register is appropriate when writing a business letter or academic paper.}
{Đây là một danh từ dùng để chỉ ngữ vực hoặc phong cách ngôn ngữ phù hợp hoàn cảnh.}

\vocabentry{Style}
{/staɪl/ (n)}
{A particular way in which something is done, created, or performed; in language, the choice of words and sentence structure.}
{The author is famous for his clear and concise writing style.}
{Đây là một danh từ dùng để chỉ văn phong hoặc phong cách.}

\vocabentry{Tone}
{/toʊn/ (n)}
{The general character or attitude of a place, piece of writing, situation, etc.}
{The tone of the email should be professional and polite.}
{Trong giao tiếp, đây là danh từ chỉ giọng điệu hoặc sắc thái tình cảm.}

\vocabentry{Emphasis}
{/ˈemfəsɪs/ (n)}
{Special importance, value, or prominence given to something.}
{The speaker put a lot of emphasis on the need for urgent environmental action.}
{Trong lời nói, đây là sự nhấn mạnh.}

\vocabentry{Clarity}
{/ˈklærəti/ (n)}
{The quality of being coherent and intelligible.}
{Lack of clarity in the instructions led to widespread confusion among the students.}
{Đây là một danh từ dùng để chỉ sự rõ ràng, trong sáng trong diễn đạt.}

\vocabentry{Conciseness}
{/kənˈsaɪsnəs/ (n)}
{The quality of being short and clear.}
{Conciseness is valued in business writing, where time is often limited.}
{Đây là một danh từ dùng để chỉ tính ngắn gọn, súc tích.}

\vocabentry{Misinterpretation}
{/ˌmɪsɪnˌtɜːrprɪˈteɪʃn/ (n)}
{The action of interpreting something wrongly.}
{Cultural differences can often lead to the misinterpretation of non-verbal cues.}
{Đây là một danh từ dùng để chỉ sự hiểu lầm hoặc diễn giải sai.}

\vocabentry{Understanding}
{/ˌʌndərˈstændɪŋ/ (n)}
{The ability to understand something; comprehension.}
{A shared language is a key foundation for mutual understanding between nations.}
{Đây là một danh từ dùng để chỉ sự thấu hiểu hoặc am hiểu.}

\vocabentry{Comprehension}
{/ˌkɑːmprɪˈhenʃn/ (n)}
{The action or capability of understanding something.}
{Listening comprehension is often the most difficult skill for new language learners to master.}
{Đây là một danh từ dùng để chỉ sự lĩnh hội hoặc hiểu biết.}

\vocabentry{Intelligibility}
{/ɪnˌtelɪdʒəˈbɪləti/ (n)}
{The quality of being possible to understand.}
{The intelligibility of his speech was affected by the poor quality of the microphone.}
{Đây là một danh từ dùng để chỉ tính dễ hiểu của lời nói hoặc văn bản.}

\vocabentry{Interpretable}
{/ɪnˈtɜːrprɪtəbl/ (adj)}
{Able to be interpreted in a particular way.}
{The data from the survey is interpretable in several different ways.}
{Đây là một tính từ dùng để chỉ cái gì đó có thể giải thích được.}

\vocabentry{Eloquent}
{/ˈeləkwənt/ (adj)}
{Fluent or persuasive in speaking or writing; clearly expressing or indicating something.}
{She gave an eloquent speech that moved everyone in the audience.}
{Đây là một tính từ dùng để chỉ khả năng diễn đạt hùng hồn, thuyết phục}

\vocabentry{Articulate}
{/ɑːrˈtɪkjələt/ (adj/v)}
{(Of a person) having or showing the ability to speak fluently and coherently.}
{He is a very articulate young man who can explain complex ideas clearly.}
{Đây là một tính từ/động từ chỉ người có khả năng diễn đạt tốt.}

\vocabentry{Persuasive}
{/pərˈsweɪsɪv/ (adj)}
{Good at persuading someone to do or believe something.}
{The lawyer gave a very persuasive argument that convinced the jury of his client's innocence.}
{Đây là một tính từ dùng để chỉ tính thuyết phục.}

\vocabentry{Coherent}
{/koʊˈhɪrənt/ (adj)}
{(Of an argument, theory, or policy) logical and consistent.}
{The witness was unable to provide a coherent account of the events.}
{Đây là một tính từ dùng để chỉ sự mạch lạc, chặt chẽ.}

\vocabentry{Mumble}
{/ˈmʌmbl/ (v)}
{Say something indistinctly and quietly, making it difficult for others to hear.}
{Please don't mumble; speak clearly so that everyone can hear you.}
{Đây là một động từ dùng để chỉ việc nói lầm bầm, không rõ lời.}

\vocabentry{Stutter}
{/ˈstʌtər/ (v/n)}
{Talk with continued involuntary repetition of sounds.}
{He has had a stutter since he was a small child, but it hasn't stopped him from becoming a teacher.}
{Đây là một động từ/danh từ dùng để chỉ sự nói lắp.}

\vocabentry{Stammer}
{/ˈstæmər/ (v/n)}
{Speak with sudden involuntary pauses and a tendency to repeat the initial letters of words.}
{She tended to stammer whenever she was nervous or excited.}
{Đây là một động từ/danh từ tương tự như stutter, chỉ sự nói cà lăm/ngập ngừng.}

\vocabentry{Whisper}
{/ˈwɪspər/ (v/n)}
{Speak very softly using one's breath without one's vocal cords.}
{They were whispering in the library so they wouldn't disturb the other students.}
{Đây là một động từ/danh từ chỉ việc nói thầm.}

\vocabentry{Shout}
{/ʃaʊt/ (v/n)}
{Utter a loud call or cry, typically as an expression of a strong emotion.}
{He had to shout to be heard over the noise of the machinery.}
{Đây là một động từ/danh từ chỉ việc hét hoặc la lớn.}

\vocabentry{Banter}
{/ˈbæntər/ (n/v)}
{The playful and friendly exchange of teasing remarks.}
{I enjoy the friendly banter between the colleagues in our office.}
{Đây là một danh từ/động từ chỉ lời nói đùa giỡn, giễu cợt thân thiện.}

\vocabentry{Gossip}
{/ˈɡɑːsɪp/ (n/v)}
{Casual or unconstrained conversation or reports about other people.}
{You shouldn't believe everything you hear; most of it is just idle gossip.}
{Đây là một danh từ/động từ chỉ chuyện ngồi lê đôi mách, tin đồn.}

\vocabentry{Chatter}
{/ˈtʃætər/ (v/n)}
{Talk informally about unimportant matters.}
{The children continued to chatter excitedly all the way home.}
{Đây là một động từ/danh từ chỉ việc nói nhảm, nói hươu nói vượn.}

\vocabentry{Dialogue}
{/ˈdaɪəlɔːɡ/ (n)}
{Conversation between two or more people as a feature of a book, play, or movie.}
{The play is full of witty and engaging dialogue between the main characters.}
{Đây là một danh từ dùng để chỉ lời thoại hoặc cuộc đối thoại.}

\vocabentry{Monologue}
{/ˈmɑːnəlɔːɡ/ (n)}
{A long speech by one actor in a play or movie.}
{The lead actor delivered a powerful monologue that left the audience in tears.}
{Đây là một danh từ dùng để chỉ lời độc thoại.}

\vocabentry{Debate}
{/dɪˈbeɪt/ (n/v)}
{A formal discussion on a particular topic in a public meeting or legislative assembly.}
{The proposed changes to the law have sparked a heated debate across the country.}
{Đây là một danh từ/động từ chỉ cuộc tranh luận.}

\vocabentry{Discussion}
{/dɪˈskʌʃn/ (n)}
{The action or process of talking about something in order to reach a decision or to exchange ideas.}
{We had a very productive discussion about the future of the project.}
{Đây là một danh từ dùng để chỉ cuộc thảo luận.}

\vocabentry{Negotiation}
{/nəˌɡoʊʃiˈeɪʃn/ (n)}
{Discussion aimed at reaching an agreement.}
{The two companies are currently in negotiations for a potential merger.}
{Đây là một danh từ dùng để chỉ sự thương lượng hoặc đàm phán.}

\vocabentry{Conflict}
{/ˈkɑːnflɪkt/ (n/v)}
{A serious disagreement or argument.}
{Effective communication skills can help to resolve conflicts in the workplace.}
{Trong giao tiếp, đây là danh từ chỉ sự xung đột hoặc mâu thuẫn.}

\vocabentry{Resolution}
{/ˌrezəˈluːʃn/ (n)}
{A firm decision to do or not to do something; the action of solving a problem.}
{The conflict was settled through a peaceful resolution reached after several meetings.}
{Trong giao tiếp là sự giải quyết mâu thuẫn.}

\vocabentry{Diplomacy}
{/dɪˈploʊməsi/ (n)}
{The profession, activity, or skill of managing international relations; skill in dealing with people.}
{Diplomacy is essential for maintaining peace and stability between nations.}
{Đây là một danh từ dùng để chỉ sự ngoại giao hoặc khéo léo trong giao thiệp.}

\vocabentry{Tact}
{/tækt/ (n)}
{Adroitness and sensitivity in dealing with others or with difficult issues.}
{You need to use a lot of tact when giving negative feedback to a colleague.}
{Đây là một danh từ dùng để chỉ sự khéo léo, tế nhị trong ứng xử.}

\vocabentry{Empathy}
{/ˈempəθi/ (n)}
{The ability to understand and share the feelings of another.}
{Showing empathy can help to build trust and rapport in a conversation.}
{Trong giao tiếp, đây là sự đồng cảm.}

\vocabentry{Rapport}
{/ræˈpɔːr/ (n)}
{A close and harmonious relationship in which the people or groups concerned understand each other's feelings or ideas.}
{She has a great rapport with her students and always makes them feel at ease.}
{Đây là một danh từ dùng để chỉ mối quan hệ hòa hợp, sự tương tác tốt.}

\vocabentry{Body language}
{/ˈbɑːdi ˌlæŋɡwɪdʒ/ (n)}
{The process of communicating non-verbally through conscious or unconscious gestures and movements.}
{Crossed arms are often seen as a sign of defensive body language.}
{Đây là một cụm danh từ dùng để chỉ ngôn ngữ cơ thể.}

\vocabentry{Facial expression}
{/ˌfeɪʃl ɪkˈspreʃn/ (n)}
{One or more motions or positions of the muscles beneath the skin of the face.}
{Her facial expression showed that she was surprised by the news.}
{Đây là một cụm danh từ dùng để chỉ biểu cảm khuôn mặt.}

\vocabentry{Eye contact}
{/ˈaɪ ˌkɑːntækt/ (n)}
{The act of looking directly into one another's eyes.}
{Maintaining good eye contact shows that you are paying attention and are interested in what the other person is saying.}
{Đây là một cụm danh từ dùng để chỉ sự giao tiếp bằng mắt.}

\vocabentry{Posture}
{/ˈpɑːstʃər/ (n)}
{A particular way of standing or sitting.}
{An upright posture can convey confidence and authority.}
{Trong giao tiếp phi ngôn ngữ, đây là danh từ chỉ tư thế.}

\vocabentry{Proximity}
{/prɑːkˈsɪməti/ (n)}
{Nearness in space, time, or relationship.}
{Standing too close to someone can make them feel uncomfortable and invade their personal space.}
{Trong giao tiếp, đây là khoảng cách giữa các cá nhân.}

\vocabentry{Haptics}
{/ˈhæptɪks/ (n)}
{The study of touch and human interaction.}
{Haptics varies greatly between cultures, with some being much more physically affectionate than others.}
{Đây là một danh từ chuyên ngành nghiên cứu về giao tiếp qua xúc giác/chạm.}

\vocabentry{Vocalics}
{/voʊˈkælɪks/ (n)}
{The study of the use of the voice to convey meaning.}
{Vocalics examines how the speed and pitch of our voice can change the meaning of our words.}
{Đây là một danh từ chuyên ngành nghiên cứu về cách sử dụng giọng nói - ngữ điệu, âm lượng.}

\vocabentry{Chronemics}
{/kroʊˈniːmɪks/ (n)}
{The study of the use of time in non-verbal communication.}
{Chronemics explores how being late or on time conveys different messages in different cultures.}
{Đây là một danh từ chuyên ngành nghiên cứu về cách sử dụng thời gian trong giao tiếp phi ngôn ngữ.}

\vocabentry{Paralanguage}
{/ˈpærəˌlæŋɡwɪdʒ/ (n)}
{The non-lexical component of communication by speech, for example intonation, pitch and speed.}
{Paralanguage can provide important clues about a speaker's emotions and attitudes.}
{Đây là một danh từ dùng để chỉ cận ngôn ngữ.}

\vocabentry{Signaling}
{/ˈsɪɡnəlɪŋ/ (n)}
{The transmission of a signal or signals.}
{Hand signals are often used by cyclists to indicate their intention to turn.}
{Đây là một danh từ dùng để chỉ việc phát tín hiệu trong giao tiếp.}

\vocabentry{Encoding}
{/ɪnˈkoʊdɪŋ/ (n)}
{The process of turning thoughts into communication.}
{Effective encoding requires choosing the right words and medium to convey your message.}
{Đây là một danh từ dùng để chỉ quá trình mã hóa thông tin.}

\vocabentry{Decoding}
{/diːˈkoʊdɪŋ/ (n)}
{The process of turning communication into thoughts.}
{Misunderstandings often happen during the decoding process when the receiver misinterprets the sender's message.}
{Đây là một danh từ dùng để chỉ quá trình giải mã thông tin.}

\vocabentry{Noise}
{/nɔɪz/ (n)}
{Any disturbance that interferes with the transmission of a message.}
{Background noise can make it very difficult to have a clear conversation over the phone.}
{Trong lý thuyết truyền thông, đây là sự nhiễu loạn.}

\vocabentry{Medium}
{/ˈmiːdiəm/ (n)}
{The agency or means by which something is accomplished, conveyed, or transferred.}
{Email is a popular medium for business communication due to its speed and convenience.}
{Đây là một danh từ dùng để chỉ phương tiện truyền tải.}

\vocabentry{Channel}
{/ˈtʃænl/ (n)}
{A medium through which information is transmitted.}
{Face-to-face conversation is considered a rich communication channel because it includes non-verbal cues.}
{Đây là một danh từ dùng để chỉ kênh truyền thông.}

\vocabentry{Context}
{/ˈkɑːntekst/ (n)}
{The circumstances that form the setting for an event, statement, or idea.}
{To understand the meaning of a word, you must consider its context within the sentence.}
{Đây là một danh từ dùng để chỉ bối cảnh.}

\vocabentry{Interpersonal}
{/ˌɪntərˈpɜːrsənl/ (adj)}
{Relating to relationships or communication between people.}
{Strong interpersonal skills are essential for success in management and sales.}
{Đây là một tính từ dùng để chỉ quan hệ/giao tiếp giữa cá nhân với nhau.}

\vocabentry{Intrapersonal}
{/ˌɪntrəˈpɜːrsənl/ (adj)}
{Occurring within the individual mind or self.}
{Intrapersonal communication involves self-reflection and internal dialogue.}
{Đây là một tính từ dùng để chỉ giao tiếp nội tâm/trong tâm trí một người.}

\vocabentry{Mass communication}
{/ˌmæs kəˌmjuːnɪˈkeɪʃn/ (n)}
{The imparting or exchanging of information on a large scale to a wide range of people.}
{Television and the internet are powerful tools of mass communication.}
{Đây là một cụm danh từ dùng để chỉ truyền thông đại chúng.}

\vocabentry{Broadcast}
{/ˈbrɔːdkæst/ (v/n)}
{Transmit (a program or some information) by radio or television.}
{The news is broadcast every hour on the national radio station.}
{Đây là một động từ/danh từ chỉ việc phát sóng.}

\vocabentry{Publish}
{/ˈpʌblɪʃ/ (v)}
{Prepare and issue for public sale.}
{She plans to publish her first novel early next year.}
{Đây là một động từ dùng để chỉ việc xuất bản.}

\vocabentry{Social media}
{/ˌsoʊʃl ˈmiːdiə/ (n)}
{Websites and applications that enable users to create and share content or to participate in social networking.}
{Social media has revolutionized the way we communicate and share information with the world.}
{Đây là một cụm danh từ dùng để chỉ mạng xã hội.}

\vocabentry{Digital communication}
{/ˌdɪdʒɪtl kəˌmjuːnɪˈkeɪʃn/ (n)}
{Any type of communication that relies on the use of technology.}
{Digital communication has made it possible to stay in touch with people from all over the world instantly.}
{Đây là một cụm danh từ dùng để chỉ truyền thông kỹ thuật số.}

\vocabentry{Connectivity}
{/ˌkɑːnekˈtɪvəti/ (n)}
{The state or quality of being connected or connective.}
{Improved internet connectivity has led to a boom in remote working and online education.}
{Đây là một danh từ dùng để chỉ khả năng kết nối.}

\vocabentry{Engagement}
{/ɪnˈɡeɪdʒmənt/ (n)}
{The action of engaging or being engaged.}
{The social media post generated a high level of engagement with thousands of likes and shares.}
{Trong truyền thông, đây là sự tương tác/gắn kết của khán giả.}

\vocabentry{Influence}
{/ˈɪnfluəns/ (n/v)}
{The capacity to have an effect on the character, development, or behavior of someone or something.}
{Public speakers use their language skills to influence the opinions of their audience.}
{Đây là một danh từ/động từ dùng để chỉ sự ảnh hưởng.}

\vocabentry{Manipulation}
{/məˌnɪpjuˈleɪʃn/ (n)}
{The action of manipulating someone or something in a clever or unscrupulous way.}
{Critics are concerned about the manipulation of information in some news reports.}
{Trong giao tiếp, đây là sự thao túng tâm lý hoặc thông tin.}

\vocabentry{Propaganda}
{/ˌprɑːpəˈɡændə/ (n)}
{Information, especially of a biased or misleading nature, used to promote a political cause or point of view.}
{The regime used propaganda to maintain control and suppress political dissent.}
{Đây là một danh từ dùng để chỉ sự tuyên truyền.}

\vocabentry{Censorship}
{/ˈsensərʃɪp/ (n)}
{The suppression or prohibition of any parts of books, films, news, etc.}
{Tight government censorship makes it difficult for journalists to report the truth.}
{Đây là một danh từ dùng để chỉ sự kiểm duyệt.}

\vocabentry{Transparency}
{/trænsˈpærənsi/ (n)}
{The condition of being transparent.}
{Transparency in government communication is essential for building trust with the public.}
{Trong giao tiếp, đây là sự minh bạch.}

\vocabentry{Authenticity}
{/ˌɔːθenˈtɪsəti/ (n)}
{The quality of being authentic.}
{Authenticity is key to building strong and lasting personal relationships.}
{Đây là một danh từ dùng để chỉ tính xác thực hoặc sự chân thành trong giao tiếp.}

\vocabentry{Credibility}
{/ˌkredəˈbɪləti/ (n)}
{The quality of being trusted and believed in.}
{The journalist's credibility was damaged after it was discovered that he had made up part of the story.}
{Đây là một danh từ dùng để chỉ sự uy tín hoặc độ tin cậy.}

\vocabentry{Ethics}
{/ˈeθɪks/ (n)}
{Moral principles that govern a person's behavior or the conducting of an activity.}
{Professional ethics require journalists to report the facts accurately and fairly.}
{Trong giao tiếp, đây là danh từ chỉ đạo đức.}

\vocabentry{Privacy}
{/ˈpraɪvəsi/ (n)}
{The state or condition of being free from being observed or disturbed by other people.}
{Data privacy is a major concern for many people who use the internet for personal communication.}
{Đây là một danh từ dùng để chỉ quyền riêng tư.}

\vocabentry{Identity}
{/aɪˈdentəti/ (n)}
{The fact of being who or what a person or thing is.}
{Your native language is a key part of your cultural identity.}
{Ngôn ngữ là một phần của bản sắc.}

\vocabentry{Culture}
{/ˈkʌltʃər/ (n)}
{The arts and other manifestations of human intellectual achievement regarded collectively.}
{Understanding a country's culture is essential for effective communication with its people.}
{Đây là một danh từ dùng để chỉ văn hóa.}

\vocabentry{Worldview}
{/ˈwɜːrldvjuː/ (n)}
{A particular philosophy of life or conception of the world.}
{Our language can influence our worldview and how we perceive reality.}
{Đây là một danh từ dùng để chỉ thế giới quan.}

\vocabentry{Assumption}
{/əˈsʌmpʃn/ (n)}
{A thing that is accepted as true or as certain to happen, without proof.}
{We should avoid making assumptions about people based on their accent or background.}
{Trong giao tiếp, đây là sự giả định/định kiến.}

\vocabentry{Stereotype}
{/ˈsteriətaɪp/ (n)}
{A widely held but fixed and oversimplified image or idea of a particular type of person or thing.}
{Stereotypes can lead to misunderstandings and prejudice in intercultural communication.}
{Đây là một danh từ dùng để chỉ định kiến hoặc khuôn mẫu.}

\vocabentry{Prejudice}
{/ˈpredʒədɪs/ (n)}
{Preconceived opinion that is not based on reason or actual experience.}
{Language barriers can sometimes fuel prejudice against foreign-born residents.}
{Đây là một danh từ dùng để chỉ thành kiến.}

\vocabentry{Tolerance}
{/ˈtɑːlərəns/ (n)}
{The ability or willingness to tolerate something.}
{A globalized world requires a high level of religious and linguistic tolerance.}
{Đây là một danh từ dùng để chỉ sự bao dung, chấp nhận sự khác biệt.}

\vocabentry{Diversity}
{/daɪˈvɜːrsəti/ (n)}
{The state of being diverse; variety.}
{Linguistic diversity should be celebrated and protected as part of our shared human heritage.}
{Đây là một danh từ dùng để chỉ sự đa dạng.}

\vocabentry{Inclusion}
{/ɪnˈkluːʒn/ (n)}
{The action or state of including or of being included within a group.}
{The policy promotes the inclusion of all languages in the national school curriculum.}
{Đây là một danh từ dùng để chỉ sự hòa nhập.}

\vocabentry{Barrier}
{/ˈbæriər/ (n)}
{A fence or other obstacle that prevents movement or access.}
{Language barriers can make it difficult for immigrants to access essential services.}
{Trong giao tiếp là rào cản ngôn ngữ.}

\vocabentry{Bridge}
{/brɪdʒ/ (v/n)}
{A structure carrying a road, path, railroad, or canal across a river; to close a gap.}
{Learning a new language can help to bridge the gap between different cultures.}
{Trong giao tiếp là động từ chỉ việc xóa nhòa khoảng cách/kết nối.}

\vocabentry{Literal}
{/ˈlɪtərəl/ (adj)}
{Taking words in their usual or most basic sense without metaphor or allegory.}
{A literal translation can sometimes miss the true meaning and nuances of the original text.}
{Đây là một tính từ dùng để chỉ nghĩa đen.}

\vocabentry{Figurative}
{/ˈfɪɡjərətɪv/ (adj)}
{Departing from a literal use of words; metaphorical.}
{The phrase "a sea of faces" is a figurative expression used to describe a large crowd.}
{Đây là một tính từ dùng để chỉ nghĩa bóng hoặc nghĩa biểu tượng.}

\vocabentry{Metaphorical}
{/ˌmetəˈfɔːrɪkl/ (adj)}
{Relating to or using metaphor.}
{The author uses metaphorical language to create a vivid and imaginative world.}
{Đây là một tính từ dùng để chỉ tính ẩn dụ.}

\vocabentry{Symbolic}
{/sɪmˈbɑːlɪk/ (adj)}
{Serving as a symbol.}
{White is symbolic of peace and purity in many Western cultures.}
{Đây là một tính từ dùng để chỉ tính biểu tượng.}

\vocabentry{Expressive}
{/ɪkˈspresɪv/ (adj)}
{Effectively conveying thought or feeling.}
{She has a very expressive face that reveals all her emotions.}
{Đây là một tính từ dùng để chỉ sự biểu cảm hoặc diễn cảm.}

\vocabentry{Communicative}
{/kəˈmjuːnɪkeɪtɪv/ (adj)}
{Willing to talk or impart information.}
{He is a very communicative person who enjoys meeting and talking to new people.}
{Đây là một tính từ dùng để chỉ người cởi mở, thích giao tiếp.}

\vocabentry{Informative}
{/ɪnˈfɔːrmətɪv/ (adj)}
{Providing useful or interesting information.}
{The documentary was very informative and taught me a lot about climate change.}
{Đây là một tính từ dùng để chỉ tính cung cấp thông tin/bổ ích.}

\vocabentry{Stimulating}
{/ˈstɪmjuleɪtɪŋ/ (adj)}
{Encouraging new ideas or enthusiasm.}
{We had a very stimulating discussion about the future of artificial intelligence.}
{Đây là một tính từ dùng để chỉ sự kích thích tư duy/hứng thú.}

\vocabentry{Challenging}
{/ˈtʃælɪndʒɪŋ/ (adj)}
{Testing one's abilities; demanding as well as stimulating.}
{Learning to write Chinese characters is a very challenging but rewarding task.}
{Học ngôn ngữ mới là một thử thách.}

\vocabentry{Rewarding}
{/rɪˈwɔːrdɪŋ/ (adj)}
{Providing satisfaction; worthwhile.}
{Traveling to a new country and speaking the local language is a very rewarding experience.}
{Học ngôn ngữ mới mang lại nhiều giá trị.}

\vocabentry{Essential}
{/ɪˈsenʃl/ (adj)}
{Absolutely necessary; extremely important.}
{Communication is an essential skill for success in any field of work.}
{Đây là một tính từ dùng để chỉ tính thiết yếu.}

\vocabentry{Fundamental}
{/ˌfʌndəˈmentl/ (adj)}
{Forming a necessary base or core; of central importance.}
{Literacy is a fundamental right that every child should have access to.}
{Đây là một tính từ dùng để chỉ tính chất nền tảng hoặc cơ bản.}

\vocabentry{Vital}
{/ˈvaɪtl/ (adj)}
{Absolutely necessary or important; essential.}
{Language is vital for the development and progress of human society.}
{Cực kỳ quan trọng đối với sự kết nối con người.}

\vocabentry{Significant}
{/sɪɡˈnɪfɪkənt/ (adj)}
{Sufficiently great or important to be worthy of attention.}
{There has been a significant increase in the use of digital communication over the last decade.}
{Đây là một tính từ dùng để chỉ sự đáng kể hoặc quan trọng.}

\vocabentry{Vibrant}
{/ˈvaɪbrənt/ (adj)}
{Full of energy and life.}
{Spanish is a vibrant language spoken by hundreds of millions of people worldwide.}
{Ngôn ngữ là một thực thể sống động.}

\vocabentry{Dynamic}
{/daɪˈnæmɪk/ (adj)}
{(Of a process or system) characterized by constant change, activity, or progress.}
{Language is dynamic and constantly evolving to reflect changes in society.}
{Đây là một tính từ dùng để chỉ tính năng động hoặc luôn thay đổi.}

\vocabentry{Evolving}
{/ɪˈvɑːlvɪŋ/ (adj)}
{Developing gradually, especially from a simple to a more complex form.}
{Our way of communicating is evolving as new technologies become available.}
{Đây là một tính từ dùng để chỉ sự đang tiến hóa/phát triển.}

\vocabentry{Universal}
{/ˌjuːnɪˈvɜːrsl/ (adj)}
{Of, affecting, or done by all people or things in the world.}
{Music is often called the universal language because it can move people from all cultures.}
{Cảm xúc là ngôn ngữ chung.}

\vocabentry{Global}
{/ˈɡloʊbl/ (adj)}
{Relating to the whole world; worldwide.}
{In today's global culture, English has become a primary tool for international cooperation.}
{Đây là một tính từ dùng để chỉ tính toàn cầu.}

\vocabentry{Traditional}
{/trəˈdɪʃənl/ (adj)}
{Existing in or as part of a tradition; long-established.}
{Traditional methods of communication, like letters, are becoming less common in the digital age.}
{Đây là một tính từ dùng để chỉ tính truyền thống.}

\vocabentry{Modern}
{/ˈmɑːdərn/ (adj)}
{Relating to the present or recent times.}
{Modern technology has made it easier than ever to communicate with people across the globe.}
{Đây là một tính từ dùng để chỉ tính hiện đại.}

\vocabentry{Sophisticated}
{/səˈfɪstɪkeɪtɪd/ (adj)}
{(Of a machine, system, or technique) developed to a high degree of complexity.}
{Human language is an incredibly sophisticated system of communication.}
{Đây là một tính từ dùng để chỉ sự tinh vi, phức tạp.}

\vocabentry{Complex}
{/ˈkɑːmpleks/ (adj)}
{Consisting of many different and connected parts.}
{The relationship between language and culture is extremely complex.}
{Đây là một tính từ dùng để chỉ tính phức tạp.}

\vocabentry{Intricate}
{/ˈɪntrɪkət/ (adj)}
{Very complicated or detailed.}
{The grammar rules of some indigenous languages are incredibly intricate.}
{Đây là một tính từ dùng để chỉ sự tỉ mỉ, chi tiết.}

\vocabentry{Fascinating}
{/ˈfæsɪneɪtɪŋ/ (adj)}
{Extremely interesting.}
{The study of the etymology of words is absolutely fascinating.}
{Đây là một tính từ dùng để chỉ sự hấp dẫn, lôi cuốn.}

\vocabentry{Haunting}
{/ˈhɔːntɪŋ/ (adj)}
{Poignant and evocative; difficult to forget.}
{The poem has a haunting beauty that stays in your mind long after you have read it.}
{Dùng để mô tả lời bài hát hoặc thơ ca truyền cảm.}

\vocabentry{Poignant}
{/ˈpɔɪnjənt/ (adj)}
{Evoking a keen sense of sadness or regret.}
{The final scene of the play was incredibly poignant and left many people in tears.}
{Đây là một tính từ dùng để chỉ sự sâu sắc, thấm thía.}

\vocabentry{Lyrical}
{/ˈlɪrɪkl/ (adj)}
{(Of literature, art, or music) expressing the writer's emotions in an imaginative and beautiful way.}
{His prose is so lyrical that it often feels like reading poetry.}
{Tính chất trữ tình.}

\vocabentry{Profound}
{/prəˈfaʊnd/ (adj)}
{(Of a state, quality, or emotion) very great or intense.}
{The philosopher's ideas had a profound impact on my thinking.}
{Đây là một tính từ dùng để chỉ sự sâu sắc, uyên thâm.}

\vocabentry{Captivating}
{/ˈkæptɪveɪtɪŋ/ (adj)}
{Capable of attracting and holding interest; charming.}
{She is a captivating storyteller who can hold an audience's attention for hours.}
{Đây là một tính từ dùng để chỉ sự quyến rũ, mê hoặc.}

